\begin{table*}[t]\centering\small
\caption{Testable open questions derived from the synthesis. Q1 and Q3 specify prospective tests. Q2 extends paired forwarding to separate its mediators; Q4 extends live detection to new workloads and deployments. Each gives a result that could support or revise the proposed explanation.}\label{tab:open-questions}
\setlength{\tabcolsep}{4pt}
\begin{tabularx}{\textwidth}{@{}P{.22\textwidth}P{.43\textwidth}Y@{}}\toprule
Question and unresolved mechanism & Controlled comparison & Discriminating result \\\midrule
\textbf{Q1. Which signals survive execution changes?} Do attacks identify recurring task structure or features tied to its implementation? & Keep task labels and observation rights fixed. Hold out compositions, coordination and client/scaffold versions separately, then jointly. Compare destination, volume/timing, sequence and recipient-visible feature families under declared frozen and adapted training. & Identify a feature family whose inference survives specified changes, or a failure tied to one changed condition. Report per-axis and joint transfer with attempted-run coverage; a pooled average cannot establish invariance. \\\addlinespace[4pt]
\textbf{Q2. How does feedback reach model traffic?} Which part is content, length or later action structure? & Fix initial instructions and task-required fields. Pair length-matched neutral padding with validity-preserving replacement of supplementary content; hold encoding and follow-up policy fixed where possible. Vary record count separately. Compare prefix, first-feedback and complete model views under tool-only and joint protection. & Additional held-out inference after feedback, with controls that exclude the competing path, supports that attribution. Report an effect interval against a prespecified meaningful-effect threshold. A nonsignificant difference alone does not rule out leakage. \\\addlinespace[4pt]
\textbf{Q3. Can cover preserve task semantics?} When should it stop at transport or execute in isolation? & On matched original tasks, compare unprotected execution, transport-terminated cover and isolated executed cover under declared resource budgets. Refit permitted attackers; measure bytes, delay, original-task completion, quota use, state changes and feedback entering the planner. & A privacy--utility operating point satisfying prespecified state and task tolerances supports the design. Resource exhaustion, changed business state or impaired completion identifies the constraint it violates. \\\addlinespace[4pt]
\textbf{Q4. Which inference survives task discovery?} Does cross-link association recover information hidden by waits and multiplexing? & Train on one concurrency/load and test new compositions, loads and gateways. Compare tool-only and associated tool/model views with observable flow identifiers; keep ground-truth task and stream mappings outside attacker inputs. Include real unrelated traffic and constant-output detector references. & Measure interval recall, false alarms per unit time and correctly localized inference together. Gains must survive changed workloads and loads at a stated false-alarm budget. \\
\bottomrule\end{tabularx}\end{table*}
